\part{kai}
%\input{./preambel/table-of-content}
%hier gehts los

\section{TikZ}
\begin{frame}
\frametitle{Umgebung}
\begin{itemize}
  \item Zum Zeichnen unter \LaTeX wird nachfolgend TikZ verwendet
  \item die Zeichenumgebung wird mit $\backslash$begin\{tikzpicture\} begonnen
  \item und mit $\backslash$end\{tikzpicture\} beendet
  \item TikZ bedeutet \textbf{T}ikZ \textbf{i}st \textbf{k}ein \textbf{Z}eichenprogramm
\end{itemize}
\end{frame}

\begin{frame}
\frametitle{Kreise}
\begin{table}[!h]
\begin{tabular}{lr}

$\backslash$begin\{tikzpicture\} & \\
\\
$\backslash$path[draw] (0,0) circle (2ex) 
&
\begin{tikzpicture}
  \path[draw] (0,0) circle (2ex);
\end{tikzpicture} 
\\  
\\
$\backslash$path[fill] (0,-2) circle (2ex) 
&
\begin{tikzpicture}
  \path[fill] (0,0) circle (2ex);
\end{tikzpicture}
\\
\\
$\backslash$path[fill=yellow, draw=red] (0,-4) circle (2ex)
&
\begin{tikzpicture}
  \path[fill=yellow,draw=red] (0,0) circle (2ex);
\end{tikzpicture}
\\
\\
$\backslash$end\{tikzpicture\} & \\
\end{tabular}
\end{table}
\end{frame}

\begin{frame}
\frametitle{Diagramme}
\begin{figure}[!h]
\input{./listings/diagram.tex}
\end{figure}
\end{frame}

\begin{frame}
\frametitle{Diagramme}
	\lstsettex
	\lstinputlisting{./listings/diagram.tex}
\end{frame}

\begin{frame}
\frametitle{Graphen}
\begin{figure}[!h]
\input{./listings/graph.tex}
\end{figure}
\end{frame}

\begin{frame}
\frametitle{Graphen}
	\lstsettex
	\lstinputlisting{./listings/graph.tex}
\end{frame}

\begin{frame}
\frametitle{Feistel}
\begin{figure}[!h]
\input{./listings/feistel.tex}
\end{figure}
\end{frame}

\begin{frame}
\frametitle{Feistel}
	\lstsettex
	\lstinputlisting{./listings/feistel.tex}
\end{frame}

\begin{frame}
\frametitle{Spirale}
\begin{figure}[!h]
\input{./listings/spirale.tex}
\end{figure}
\end{frame}

\begin{frame}
\frametitle{Spirale}
	\lstsettex
	\lstinputlisting{./listings/spirale.tex}
\end{frame}

